\section{Research Methods for Interactive Systems Design}

\begin{ejercicio}
    Briefly describe the design space of an interactive system.

    The design space of an interactive system refers to the range of possible designs for that system. Three main aspects need to be taken into account:
    \begin{itemize}
        \item \ul{User}: The characteristics, needs, and preferences of the users who will interact with the system.
        \item \ul{Situation}: The current state of the problem, how is it currently solved, and what are the limitations of the current solution.
        \item \ul{Context}: The environment in which the system will be used, including the physical, social, and cultural context.
    \end{itemize}
\end{ejercicio}

\begin{ejercicio}
    What is the Hawthorne effect? How can it affect the results of a research study? Provide an example of a research study where the Hawthorne effect could be present and explain how you would mitigate it.

    The Hawthorne effect refers to the phenomenon where individuals modify their behavior in response to being observed or studied. This can affect the results of a research study by introducing bias, as participants may not behave naturally, leading to inaccurate conclusions. A research method that could be affected by the Hawthorne effect is Ethnography, where the researcher observes participants in their natural environment. Some examples of research studies where the Hawthorne effect could be present include:
    \begin{itemize}
        \item A study on the effectiveness of a new teaching method in a classroom setting, where the researcher is observing the students and teachers. This could be mitigated by using unobtrusive observation techniques, such as video recording or using a one-way mirror, to minimize the impact of the researcher's presence on the participants' behavior.
    \end{itemize}
\end{ejercicio}

\begin{ejercicio}
    Explain the concept of ``going native'' in ethnography. To which type of ethnography does it refer? How can a researcher mitigate the risk of going native?

    ``Going native'' in ethnography refers to the process where the researcher becomes so immersed in the culture or community they are studying that they lose their objectivity and begin to adopt the behaviors, beliefs, and values of the participants. This can lead to biased observations and interpretations, as the researcher may no longer be able to maintain a critical distance from the subject of study. It is particularly relevant to ``Participating Observation'' ethnography, where the researcher actively engages with the community and participates in their daily activities. To mitigate the risk of going native, researchers can maintain a reflective journal, seek feedback from peers, and regularly step back to analyze their observations from an objective standpoint.
\end{ejercicio}

\begin{ejercicio}
    Briefly describe how ethnography works. Highlight its advantages and limitations, and give an example application in a real-world scenario where it would be particularly useful.

    Ethnography is a qualitative research method that involves the systematic study of people and cultures in their natural environment. Researchers immerse themselves in the community they are studying, often for extended periods, to observe and interact with participants. The goal is to gain a deep understanding of the social dynamics, behaviors, and cultural practices of the group.

    Advantages of ethnography include its ability to provide rich, detailed insights into human behavior and social interactions, as well as its flexibility in adapting to the context of the study. However, it also has limitations, such as the potential for researcher bias, the time-consuming nature of data collection, and challenges in generalizing findings to larger populations.

    An example application of ethnography could be studying how the information flows within a hospital setting to improve patient care. By observing the interactions between healthcare professionals, patients, and administrative staff, researchers can identify bottlenecks, communication breakdowns, and areas for improvement in the hospital's processes.
\end{ejercicio}

\begin{ejercicio}
    Compare the roles of the researcher in Ethnography, Action Research, and Research through Design. What are the similarities and differences?
    \begin{itemize}
        \item In Ethnography, the researcher takes on the role of an observer and participant, immersing themselves in the community or culture they are studying to gain a deep understanding of social dynamics and behaviors. The researcher aims to remain as objective as possible while collecting qualitative data through observations, interviews, and participation.

        \item In Action Research, the researcher is an active participant in the problem-solving process, collaborating with stakeholders to identify issues, implement interventions, and evaluate outcomes. The researcher is both a facilitator and a learner, working closely with participants to co-create solutions and improve practices.

        \item Lastly, in Research through Design, the researcher takes on the role of a designer and innovator, using design processes to explore and generate knowledge. The researcher creates prototypes, tests them, and iteratively refines their designs based on feedback and observations. Here, no inmersion in a community is required, but the researcher is actively involved in the design process.
    \end{itemize}

    In all of the methods, the researcher is actively engaged in the research process, and he tries to get information about the design space of the problem. However, the level of immersion and the specific role of the researcher vary across the methods. Ethnography emphasizes observation and understanding, Action Research focuses on collaboration and problem-solving, and Research through Design centers on innovation and iterative design.
\end{ejercicio}

\begin{ejercicio}
    Explain the concept of ``Refreezing'' in Action Research. What other two stages are part of the Action Research cycle? How do they relate to each other?

    Action Research is a cyclical process that involves three main stages: Planning, Action, and Evaluation.
    \begin{itemize}
        \item \ul{Planning}: In this stage, the researcher identifies the problem, gathers relevant information, and develops a plan of action to address the issue. Here, ``unfreezing'' occurs, as the researcher and participants recognize the need for change and prepare to challenge existing practices or beliefs.
        \item \ul{Action}: In this stage, the researcher and participants implement the planned interventions or changes. They actively engage in the process of ``change'', testing new approaches and strategies to address the identified problem.
        \item \ul{Evaluation}: In this stage, the researcher and participants assess the outcomes of the implemented changes, reflecting on their effectiveness and identifying areas for further improvement. This stage is where ``refreezing'' occurs, as the new practices or beliefs are solidified and integrated into the existing system or culture, ensuring that the changes are sustained over time.
    \end{itemize}
\end{ejercicio}

\begin{ejercicio}
    State the advantages and disadvantages of two of the research methods that were presented in the lecture. In this context, explain why combining methods might be beneficial. What is this process called?

    Ethnography has the advantage of providing rich, detailed insights into human behavior and social interactions, allowing researchers to gain a deep understanding of the context in which people operate. However, it can be time-consuming, resource-intensive, and difficult to generalize findings to larger populations. Survey research, on the other hand, allows for the collection of data from a large number of participants relatively quickly and efficiently and helps deciding if results can be generalized to a larger population. However, it may not provide the reasoning behind participants' responses and can be limited by the quality of the survey design.

    Triangulation is the process of combining multiple research methods to gain a more comprehensive understanding of a research problem. By using different methods, researchers can cross-validate findings, reduce bias, and increase the reliability and validity of their results. For example, combining ethnography with survey research can provide both in-depth qualitative insights and quantitative data, allowing for a more holistic understanding of the research problem.
\end{ejercicio}

\begin{ejercicio}
    Imagine you are developing a platform for interaction between different people. The platform concept is ready and you want to gain user feedback. Which research method(s) do you use and why?

    Since the platform concept is ready and user feedback is desired, field deployment would be an appropriate research method to use. Field deployment involves releasing the platform to a select group of users in their natural environment, allowing them to interact with the system and provide feedback based on their experiences. This method allows for real-world testing and evaluation of the platform's usability, functionality, and overall user experience. Additionally, user interviews or surveys can be conducted alongside field deployment to gather qualitative and quantitative feedback from participants, providing a comprehensive understanding of user needs and preferences.
\end{ejercicio}

\begin{ejercicio}
    A new social robot is developed. It is supposed to help seniors in an elderly home and also support them with being more active in their daily life. Which research methods should be used to evaluate the acceptance of the robot and its impact on the elderly? Describe the research design.

    Triangulation of research methods would be appropriate to evaluate the acceptance of the social robot and its impact on the elderly. On the one hand, a survey research method can be used to gather quantitative data on the seniors' acceptance of the robot, as there are standardized questionnaires available to measure technology acceptance. On the other hand, ethnography can be used to gather qualitative data on the seniors' experiences and interactions with the robot in their natural environment, understanding also the social dynamics and cultural context of the elderly home.

    The population of the study would consist of seniors living in an eldery home, but the sampling frame and the sample would be only those living in the elderly home where the robot is deployed. The sample size would depend on the number of seniors living in the elderly home, but a minimum of 30 participants would be recommended to ensure sufficient data for analysis.
\end{ejercicio}

\begin{ejercicio}
    You are an employee in a company, and your boss has assigned you the task of gathering research data for a new program that will be used in schools. Select one of the research methods and explain why this method is appropriate for this scenario.

    Since no direct action is needed, but rather the goal is to gather research data, ethnography would be an appropriate research method to use in this scenario. Ethnography allows for the systematic study of people and cultures in their natural environment, providing rich, detailed insights into the social dynamics, behaviors, and cultural practices of the group being studied. By immersing oneself in the school environment, the researcher can observe and interact with students, teachers, and staff to gain a deep understanding of their needs, preferences, and challenges related to the new program. This information can then be used to inform the design and development of the program to ensure it meets the needs of its intended users.
\end{ejercicio}

\begin{ejercicio}
    You want to gather knowledge about how people in an office communicate with each other and distribute different tasks between each worker. Later on, you want to design a new technical communication tool to make communication and task distribution easier. Which research method would you use to gather insights and why?

    Ethnography would be an appropriate research method to gather insights about how people in an office communicate and distribute tasks. By immersing oneself in the office environment, the researcher can observe and interact with employees to gain a deep understanding of their communication patterns, task distribution processes, and any challenges they face. This qualitative data can provide valuable insights into the needs and preferences of the users, which can then inform the design of a new technical communication tool that addresses their specific requirements.
\end{ejercicio}

\begin{ejercicio}
    Please indicate for the following cases below which study method you would try to use to solve the following issue and why, as well as how many participants (approximately) you would recruit and why. Consider everything you have learned in the lecture and teaching cases.
    \begin{enumerate}
        \item You want to design a new HoloLens application that is supposed to help care workers during their daily activities in intensive care units. However, your knowledge about the domain and working conditions is very limited. So, you decide to gather some knowledge about intensive care units first. How would you want to gather knowledge to design your application?
        
        Ethnography would be an appropriate research method to gather knowledge about intensive care units. By immersing oneself in the ICU environment, the researcher can observe and interact with care workers to gain a deep understanding of their daily activities, challenges, and needs. Since the research would not have the needed knowledge about the domain, it would be observing participation rather than participating observation. The number of participants would depend on the size of the ICU and the number of care workers available, but a minimum of 10-15 participants would be recommended to ensure sufficient data for analysis.
    \end{enumerate}
\end{ejercicio}
